\documentclass[a4paper,12pt]{article}

% 基础宏包配置
\usepackage{ctex}      % 中文支持
\usepackage{geometry}  % 页面布局
\usepackage{titlesec}  % 标题格式
\usepackage{fancyhdr}  % 页眉页脚
\usepackage{xcolor}    % 颜色支持
\usepackage{setspace}  % 行间距

% 页面边距设置 - 稍微收紧边距以增强文本墙效果
\geometry{left=2.5cm, right=2.5cm, top=2.5cm, bottom=2.5cm}
\onehalfspacing 

% 颜色定义
\definecolor{primaryBlue}{RGB}{0, 56, 107}

% 标题格式定制
\titleformat{\section}
{\normalfont\Large\bfseries\color{primaryBlue}}{\thesection}{1em}{}

% 页眉页脚
\pagestyle{fancy}
\fancyhf{}
\lhead{\small \textbf{2024年度ESG报告}}
\rhead{\small 迈向零碳未来}
\cfoot{\thepage}

\title{\textbf{\huge 2024年度环境、社会及治理（ESG）报告}\\ \vspace{0.5em} \large 科技驱动与责任同行}
\author{本公司董事会办公室}
\date{2024年3月}

\begin{document}

\maketitle

\section{战略篇：新质生产力与稳健治理}

2024年是全球汽车产业转型的关键之年，在国家“双碳”战略与“新质生产力”宏观政策的指引下，本公司坚定不移地推进电动化、智能网联化、高端化、国际化的“三直四化”战略，致力于构建绿色、高效的出行生态。在这一年里，公司克服了复杂的外部环境挑战，实现了经营业绩的稳步增长与治理水平的持续提升。

在经济绩效与股东回报方面，得益于海外市场的拓展与高端产品的热销，公司全年实现营业收入372.18亿元，同比增长达到37.63\%，归属于母公司股东的净利润为41.16亿元。为了切实回馈股东的信任，践行长期主义价值，公司在报告期内实施了两次现金分红，全年累计发放现金股利高达44.28亿元，分红总额超过了当期净利润，充分彰显了公司稳健的现金流与回馈投资者的诚意。与此同时，我们深知合规经营是企业发展的生命线，公司已连续第13年获得交易所信息披露工作“A级”评价，并顺利通过了ISO 37301合规管理体系认证，标志着内部治理达到了国际标准。在供应链管理上，我们坚持构建共生共赢的生态圈，目前共有合作供应商530家，其中20\%来自河南本地，有效带动了区域产业集群发展；为营造廉洁透明的商业环境，我们已与超过900家供应商（含潜在供应商）签署了廉洁协议，严守商业道德底线。

\section{社会篇：创新驱动与人文关怀}

科技创新是推动可持续发展的核心动力，而对人的关怀则是企业社会责任的温度所在。2024年，公司坚持高强度的研发投入，全年研发支出达17.88亿元，占营业收入的4.80\%，汇聚了包括16名博士在内的3,505名研发人员。这一投入转化为了丰硕的知识产权成果，全年新增授权专利176项，其中发明专利114项，累计参与制定标准317项。

为了让公众更好地理解我们的技术成果，我们将核心技术进行了通俗化解读。我们在行业内率先推出了“YOS车机操作系统”，该系统实现了软件与硬件的解耦，就像智能手机一样，车辆可以通过OTA空中下载技术进行远程升级，实现了功能的常用常新；同时，我们研发的“C架构”通过中央计算与区域控制的模式，如同给汽车装上了超级大脑，大幅减少了线束使用并提升了响应速度；在核心零部件方面，我们的多合一动力域控制器应用了碳化硅模块，相比传统硅基模块，开关损耗降低了50\%至70\%，显著提升了能源转化效率；针对电动车冬季续航难题，我们开发了多源超低温热泵技术，能从空气和余热中提取热量，实测采暖节能效果达30\%。在智能网联领域，公司的L4级自动驾驶巴士已累计安全运营超过1,700万公里，并获批国家首批智能网联汽车准入试点。

在员工权益保护方面，截至报告期末，集团员工总数为16,707人，其中女性占比10.02\%。我们始终将员工安全放在首位，全年投入劳动防护用品资金5,500余万元，环境改善资金2,200余万元，以及安全培训资金180余万元。通过对7.9万人次的安全教育培训，我们实现了全年无新增职业病的优异绩效。在社会公益方面，公司全年对外捐赠及公益总投入超4,000万元。其中，“儿童交通安全公益行”项目足迹遍布13个省市，我们在华东某市的活动中通过模拟大车盲区，让近万名儿童直观体验了“看不见的危险”；在海外，我们在南美某国发起的“零碳森林”项目已累计植树36,000余棵，既改善了当地生态，也传递了中国企业的绿色承诺。

\section{环境篇：低碳运营与绿色制造}

面对全球气候变化挑战，我们承诺将绿色理念贯穿于生产运营的全生命周期。2024年，通过精细化管理，公司的碳排放强度控制在0.07吨二氧化碳/万元营收。具体排放数据方面，范围1直接排放量为57,330.53吨二氧化碳，范围2外购电力排放量为122,786.63吨二氧化碳，综合能源消费量为34,560.94吨标准煤。

为了降低环境足迹，我们实施了一系列务实的节能减排行动。全年共实施节能改造项目7项，通过设备升级与工艺优化，年节电1,258兆瓦时，年节气11.34万立方米；同时，我们建立了完善的余热回收系统，全年回收利用热能约29,306吉焦，大幅减少了化石能源消耗。在能源结构转型上，公司积极引入清洁能源，自有光伏发电量达155.68万千瓦时，累计减少碳排放约835吨，并额外采购绿色电力13,580兆瓦时。在资源循环利用方面，我们坚持“一水多用”，全年取水量141.25万吨，重复用水量高达6,465.36万吨，重复用水率达到97.8\%。废弃物管理同样严格合规，一般工业固废中可回收利用量为61,485.76吨，危险废物产生量4,110.84吨并实现了100\%合规处置。此外，我们在物流环节大力推行绿色包装，国产零部件的周转包装占比已达到89\%，有效减少了包装废弃物的产生。未来，我们将继续在氢能技术、供应链碳管理等领域探索，为实现碳中和目标贡献力量。

\vspace{3cm}

\end{document}